\begin{frame}{왜 마진은 $\frac{2}{\|w\|}$인가 (1/2)}
  $w$ 벡터는 결정 경계에 \textbf{수직}이다 —
  $w^\top x + b = c$인 모든 점이 $w$와 수직인 초평면에 있기 때문.
  \begin{itemize}
    \item ``경계에서 어떤 점까지의 최단 거리''는
      $w$의 방향으로만 측정하면 된다.
    \item $w$ 방향으로 1만큼(유클리드 거리) 움직이면
      $w^\top x + b$의 값은 얼마나 변할까?
    \item $w$ 방향의 단위벡터 $\hat w = w/\|w\|$,
      $w^\top x$는 $w$ 방향으로 이동할 때 \textbf{선형}으로 변하므로
      단위 이동마다 변하는 양 = 정확히 $\|w\|$.
    \item 역으로 $w^\top x + b$가 1만큼 변하려면
      $w$ 방향으로 $\frac{1}{\|w\|}$만 이동하면 된다.
  \end{itemize}
\end{frame}
